% Sections 5.2 and 5.3 merged: the two controls that answer the two objections
% to the headline, stated with their results rather than with their
% construction. Design details live in 5.1's table caption and the appendix.

\subsection{Corpus size and region selection}
\label{sec:size-controls}

Two objections do not require argument, because both quantities can be varied
directly. The first is arithmetic: more proofs find more distinct premises
whatever the populations are doing, so a vocabulary compared at unequal $n$
confounds breadth with sample size. We rarefy the human corpus (seeded draws
of $n$ proofs, distinct premises counted, mean reported), which gives its
vocabulary at every corpus size. At the machine corpus's own size,
\machineProofs{} proofs, the human vocabulary is \humanVocabRarefied{}
premises, so equalizing size moves the ratio from \vocabRatio{}$\times$ to
\vocabRatioRarefied{}$\times$ rather than to one. Further down the curve the
point is sharper still. At a twentieth of the human corpus,
\rarefiedFivePctProofs{} proofs and far fewer than the machine corpus has, the
human vocabulary is \humanVocabRarefiedFivePct{} premises, an order of
magnitude above what the machine corpus reaches at full size. Neither curve has
saturated, so both reported vocabularies are lower bounds on their
populations', and the machine curve is the nearer of the two to flattening; we
do not extrapolate either beyond the corpora we have.

The second objection is the genre difference named in
Section~\ref{sec:background}: if competition problems occupy a small corner of
the library, a narrow machine vocabulary is a fact about which problems were
posed. The control we would prefer, a same-problem comparison, does not exist.
The only human Lean~4 proofs of miniF2F statements are a handful released as a
prover's own few-shot prompt pool, the LeanDojo Lean~4 port carries
\texttt{sorry} throughout, and the original corpus is Lean~3. We therefore
restrict by region instead, keeping every top-level region that contains at
least one machine premise. The rule was fixed before the run, is support-based so there
is no threshold or top-$k$ to tune, and introduces no metric. The machine
vocabulary occupies \areaRegionsMachine{} of the \areaRegionsHuman{} regions
the human vocabulary spans, and those regions hold \areaHumanSurvivalPct{}\% of
human premises, so the restriction discards most of the human vocabulary rather
than trimming it. Inside them the ratio is \areaRestrictedRatio{}$\times$ at a
Jaccard of \areaRestrictedJaccard{}, with \areaRestrictedHumanOnlyPct{}\% of
the restricted human vocabulary still untouched.

Machine leakage is zero by construction under this rule, which is what makes
it primary and its registered sensitivities the check on it. The
top-twelve-namespace variant gives \areaTopTwelveRatio{}$\times$ but discards
\areaTopTwelveMachineLeakPct{}\% of the machine vocabulary on the way, and a
restriction that drops nearly a third of the population it controls for is
answering a different question. We report the range
\areaTopTwelveRatio{}$\times$ to \areaRestrictedRatio{}$\times$ rather than
either endpoint. Neither control reaches genre, or selection \emph{within} a
region. \texttt{Nat} holds competition arithmetic and analytic number theory
alike. The closest genre check we have is exploratory, and we state it here rather than
only in limitations because it is the design that most constrains what the
headline can mean. Mathlib's own IMO archive, against the machine corpus
rarefied to the same proof count, gives \genreCheckRatio{}$\times$, the
smallest gap of any design we ran. It was not pre-registered, it is a lower
bound, it is not extraction-audited, and it is not regenerable from the
released artifact (Section~\ref{sec:limitations}, L6). These designs bound how much of the gap corpus size and region selection
explain; Section~\ref{sec:limitations}, L4, bounds what the result can mean in
the presence of what they do not.
